%%%%%%%%%%%%%%%%%%%%%%%%%%%%%%%%%%%%%%%%%
% Beamer Presentation
% Ricks Template
%
%
%%%%%%%%%%%%%%%%%%%%%%%%%%%%%%%%%%%%%%%%%

%----------------------------------------------------------------------------------------
%	PACKAGES AND THEMES
%----------------------------------------------------------------------------------------
\documentclass[t,aspectratio=169,11pt,presentation]{beamer}
%\documentclass[t,aspectratio=169,11pt,presentation]{beamer}
%\documentclass[t,aspectratio=169,11pt,draft]{beamer}


% Define colors
\definecolor{ptb5}{RGB}{54,75,154}	 % NT                   
\definecolor{ptb4}{RGB}{74,123,183} % RC + NT	                    
\definecolor{ptb3}{RGB}{110,166,205}                     
\definecolor{ptb2}{RGB}{152,202,225}% EC + RC + NT                  
\definecolor{ptb1}{RGB}{194,228,239}
\definecolor{ptg0}{RGB}{234,236,204}
\definecolor{ptr1}{RGB}{254,218,139} %EC only
\definecolor{ptr2}{RGB}{253,179,102}	                    
\definecolor{ptr3}{RGB}{246,126,75} % AT +	EC + RC                    
\definecolor{ptr4}{RGB}{221,61,45}	 % AT + EC                   
\definecolor{ptr5}{RGB}{165,0,38}	 % AT	   

\usetheme{Boadilla}




% Footer and page numbers

% First remove the navigation symbols from the bottom of all slides 
\setbeamertemplate{navigation symbols}{} 

% Remove the colored line at the bottom
\setbeamertemplate{footline}{}


% Customize items
\setbeamertemplate{itemize items}{\color{ptb2}{\rule[0.05em]{0.6em}{0.6em}}}
\setbeamertemplate{itemize subitem}{\color{ptb2}{\raisebox{0.12em}{$\blacktriangleright$}}}
\setbeamercolor{enumerate item}{fg=black}
\setbeamertemplate{enumerate item}[default]
\setbeamercolor{enumerate subitem}{fg=black}
\setbeamertemplate{enumerate subitem}[default]


\setbeamercolor{subsection in toc}{bg=white,fg=black}




% Table of Contents (if we have one)
\setbeamertemplate{section in toc}{\inserttocsectionnumber.~\inserttocsection}
\setbeamercolor{section in toc}{bg=white,fg=black}
\setbeamertemplate{subsection in toc}{%
  \hspace{1.2em}{{\color{ptb5}{\rule[0.05em]{0.6em}{0.6em}}}~\inserttocsubsection\par}
}
\setbeamercolor{subsection in toc}{bg=white,fg=black}


% Customize beamer colors

\setbeamercolor{titlelike}{fg=white,bg=ptb5}
\setbeamercolor{button}{bg=ptb4,fg=white}
%}

% Nicer looking Itemize and enumerate environments
\newenvironment{wideitemize}{\itemize\addtolength{\itemsep}{14pt}}{\enditemize}
\newenvironment{wideenumerate}{\enumerate\addtolength{\itemsep}{14pt}}{\endenumerate}

% Import packages
\usepackage[default]{lato} % More modern looking font than Helvetica
\usepackage{geometry,caption,color,setspace,dsfont,commath,amsmath,amssymb,bm,mathrsfs}
\usepackage[normalem]{ulem}
\usepackage[comma]{natbib}
\usepackage[short]{optidef}
\usepackage{hhline}
\usepackage{array}
\usepackage{booktabs} % Allows the use of \toprule, \midrule and \bottomrule in tables
\usepackage{adjustbox}
\usepackage{graphicx,xcolor,bbm,xcomment}
\usepackage{appendixnumberbeamer}
\usepackage{textcomp}
\usepackage{colortbl}
\usepackage{subcaption} 
\usepackage{tikz}
\usetikzlibrary{calc,patterns,positioning}
\usepackage{pdflscape}
\usepackage{environ}
\usepackage{array}
\usepackage{etoolbox}
\usepackage[normalem]{ulem}
\BeforeBeginEnvironment{itemize}{\par}
\BeforeBeginEnvironment{enumerate}{\par}

\bibliographystyle{ecta}

\newcommand*\diff{\mathop{}\!\mathrm{d}}


% Make Section Header Frames
\AtBeginSection[]{

  \setbeamercolor{titlelike}{bg=ptb5,fg=white}
  \setbeamertemplate{navigation symbols}{} 
  \begin{frame}[noframenumbering]
  \vfill
  \centering
  \begin{beamercolorbox}[sep=8pt,center,shadow=true,rounded=true]{title}
    \usebeamerfont{title}\insertsectionhead\par%
  \end{beamercolorbox}
  \vfill
  \end{frame}
  % Make a new prettier page number
\addtobeamertemplate{navigation symbols}{}{%
    \usebeamerfont{footline}%
    \hspace{5em}%
    {\color{black!50}{\small\insertframenumber/\inserttotalframenumber}}%
}
\setbeamercolor{titlelike}{fg=ptb5,bg=white}
}

\AtBeginSubsection[]{
  \setbeamertemplate{navigation symbols}{} 
  \begin{frame}[noframenumbering]
  \vfill
  \centering
  \begin{beamercolorbox}[sep=8pt,center,shadow=false,rounded=true]{title}
    \usebeamerfont{title} \insertsubsectionhead\par%
  \end{beamercolorbox}
  \vfill
  \end{frame}
  % Make a new prettier page number
\addtobeamertemplate{navigation symbols}{}{%
    \usebeamerfont{footline}%
    \hspace{5em}%
    {\color{black!50}{\small\insertframenumber/\inserttotalframenumber}}%
}
}



\tikzset{
    invisible/.style={opacity=0},
    visible on/.style={alt={#1{}{invisible}}},
    alt/.code args={<#1>#2#3}{%
      \alt<#1>{\pgfkeysalso{#2}}{\pgfkeysalso{#3}} % \pgfkeysalso doesn't change the path
    },
  }
  
  
\makeatletter
\newsavebox{\measure@tikzpicture}
\NewEnviron{scaletikzpicturetowidth}[1]{%
  \def\tikz@width{#1}%
  \def\tikzscale{1}\begin{lrbox}{\measure@tikzpicture}%
  \BODY
  \end{lrbox}%
  \pgfmathparse{#1/\wd\measure@tikzpicture}%
  \edef\tikzscale{\pgfmathresult}%
  \BODY
}
\makeatother

\tikzset{
  pics/student/.style args={#1,#2}{
     code={
       \filldraw [fill=#1,draw=black,visible on=#2] (.05,.15) .. controls (.05,.7) and (.65,.8) .. (.75,.2) .. controls (.55,.05) and (.05,0) .. (.05,.15) ;
       \filldraw [fill=#1,draw=black,visible on=#2] (.35,.75) circle [radius=.22];
     }
  }
}

%----------------------------------------------------------------------------------------
%	TITLE PAGE
%----------------------------------------------------------------------------------------
 
\title{Are Value Added Measures Regressive? Evidence from Applying Welfare Theory} 

\author{Tanner Eastmond \hfill \and Nathan Mather    \\   \textit{teastmond@ucsd.edu}\hfill \textit{njmather{\fontfamily{qag}\selectfont @}umich.edu} \\ {\color{white}{t}} \\  Michael D Ricks  \and \hfill  Julian Betts \\
\textit{ricksmi{\fontfamily{qag}\selectfont @}umich.edu} \hfill \textit{jbetts@ucsd.edu}} 
\institute{}
\date{\today} 
% Title paage ________________________________________________________________

\begin{document}

\begin{frame}[noframenumbering]
\titlepage 
\end{frame}


\setbeamercolor{titlelike}{fg=ptb5,bg=white}
% Make a new prettier page number
\addtobeamertemplate{navigation symbols}{}{%
    \usebeamerfont{footline}%
    \hspace{5em}%
    {\color{black!50}{\small\insertframenumber/\inserttotalframenumber}}%
}

%------------------------------------------------
% What is my question and how will I answer it?
\begin{frame}{\textbf{Teacher value added measures ``mean'' well}}

\begin{wideitemize}
    \item 30+ states use value added measures to rank, evaluate, or compensate teachers
    
    \item<2-> Value added and similar mean-like quality measures relevant to policy
    
    \item <3->{\color<2-4>{ptr3}\textbf{Means}} may not capture all of the socially relevant information
        %  spend \$50B on 4M kindergarteners each year
        % In most cases this is a question of when not a question of whether
    \begin{itemize}
        \item<4-> {\color<4>{ptr5}\textbf{Impacts:}} heterogeneity means comparative advantage shapes overall {\color<4>{ptb4}\textbf{gains}}
        \item<5-> {\color<5>{ptr5}\textbf{Preferences:}} heterogeneity means policymakers care about {\color<5>{ptr2}\textbf{gaps}}
    \end{itemize}
    
    \item<6-> {\color<6>{ptr5}\textbf{Positive: Does value added vary meaningfully over student achievement?}} 
  
    \item<7> {\color<7>{ptr5}\textbf{Normative: How does ignoring heterogeneity affect equity and efficiency?}}
  
\end{wideitemize}
\end{frame}

%------------------------------------------------
% What is my question and how will I answer it?
\begin{frame}{\textbf{If there is heterogeneity, how much does it matter?}}

\begin{wideitemize}
    \item \textbf{Standard VA:} Do traditional measures of value added capture real effects?
    
    \begin{itemize}
        \item<2-> Yes! Teacher effects on student achievement seem to be unbiased
        
        {\tiny \color{gray}\citep[e.g.,][etc.]{chetty2014measuring1,Koedel2015}}
        
        \item<3-> High value-added teachers cause improvements in long-term student outcomes
        
        {\tiny \color{gray}\citep[e.g.,][]{chetty2014measuring2,pope2017multidimensional}}
        
        
    \end{itemize}
    
    \item<4-> \textbf{Heterogeneity:} Is there heterogeneity in each teacher's value added?
    
    \begin{itemize}
        \item<5-> Yes? Emerging evidence suggests differences across student types
    
        {\tiny \color{gray} \citep[e.g.,][]{lockwood2009,condie2014teacher,pope2017multidimensional,Delgado2020,bates2022teacher}}
    
    \end{itemize}

    \item<6-> Does heterogeneity matter beyond {\color{ptr3}\textbf{mean}} estimates (e.g., for long-term effects)? 
        
    \item<7-> Heterogeneous estimates can't rank teachers or compare (non-Pareto) allocations... 
    
    
    %cannot inform counterfactual policy discussions about reducing noncompliance or economic questions about efficiency and equity. This is because these questions hinge on the sorting of and effects on noncompliant children---patterns the complier-based LATE cannot inform.
    

\end{wideitemize}
\end{frame}

%------------------------------------------------
% What is my question and how will I answer it?
\begin{frame}{\textbf{This paper evaluates heterogeneous value added}}

\begin{wideitemize}
    \item Combine models of heterogeneous value added with welfare theory
    \begin{itemize}
        \item Explore theoretical justification and practical implications of heterogeneity
        \item Propose two estimators: two groups and nonparametric
        \item Aggregate to welfare using {\color<2>{ptr5}\textbf{explicit}} welfare weights
    \end{itemize}
    \item<2-> Estimate value added using data from the San Diego Unified School District
    \begin{itemize}
        \item Test score effects by achievement
        \item Long-term effects of value added vs heterogeneous value added
    \end{itemize}
    
    \item<3-> Characterize implications for welfare and policy
    \begin{itemize}
        \item Use welfare weights to solve reallocation as  mixed-integer linear programming problem
        \item Compare implications of different measures for students and teachers
    \end{itemize}


\end{wideitemize}
\end{frame}

%------------------------------------------------
% What is the answer
\begin{frame}{\textbf{Ignoring heterogeneity curtails scores and fosters inequity}}

            
\begin{wideitemize}
%\item There is significant selection:
% Recall I am thinking about selection in terms of potential untreated outcomes
\item {\color<1>{ptr5}\textbf{Heterogeneity:}} Replicate and extend results that teachers have different effects

{\tiny \color{gray}\citep{lockwood2009,condie2014teacher,bates2022teacher}}

\begin{itemize}
    \item<2-> {\color<2>{ptr5}\textbf{Long Term:}} Value Added ``Match'' raises graduation and college enrollment beyond VA
\end{itemize}


\item<3-> {\color<3>{ptr5}\textbf{Student Welfare:}} Reallocations can raise scores and shrink gaps 

%{\tiny \color{gray} \citep{bedard2006persistence,black2011too,mccrary2011effect,cook2016birthdays,molnar2018redshirting,fortner2019forced}}

\begin{itemize}
    \item  Raise math scores by 0.09 student standard deviations (69\% beyond standard VA)
    
    \item  Shrink racial math gap by 0.12 student standard deviations (standard VA widens gap)
    
    %{\tiny \color{gray} \citep[]{lalonde1986evaluating,brand2010benefits,einav2010optimal,walters2018demand,finkelstein2019take,ito2021selection}}

    
    %{\tiny \color{gray} \citep[]{lalonde1986evaluating,brand2010benefits,einav2010optimal,walters2018demand,finkelstein2019take,ito2021selection}}
    
\end{itemize}

%brinch2017beyond,mogstad2018using,kline2019heckits,walters2018demand,ito2021selection


\item<4-> {\color<4>{ptr5}\textbf{Teacher Welfare:}} Standard VA methods penalize minority teachers

\begin{itemize}
    \item Ignoring heterogeneity would be a 7\% performance pay tax on minorities
    
    
%{\tiny \color{gray} \citep[]{graue2000redshirting,deming2008lengthening,illinois2019anact}}
\end{itemize}





\end{wideitemize}

\end{frame}





\setbeamertemplate{navigation symbols}{} 
  
%-----------------------------------------------
\begin{frame}[c,noframenumbering]{\textbf{Today's Talk}}

\begin{columns}
\begin{column}{0.575\textwidth}
\begin{wideenumerate}
    \item How Do We Model and Measure Value Added?
    \begin{itemize}
        \item Conceptual model
        \item Empirical methods
    \end{itemize}
   
    \item How Do Teachers Affect Different Students?
    \begin{itemize}
        \item Test-score impacts
        \item Long-term impacts
    \end{itemize}

    \item What About Welfare and Policy?
    \begin{itemize}
        \item Comparing reallocations
        \item Recovering implicit weights
        \item Reranking teachers
    \end{itemize}
\end{wideenumerate}
\end{column}
\begin{column}{0.40\textwidth}
\vspace{-1.5em}
\begin{tikzpicture}
%\draw[line width=.5cm] (0,0)  -- (3.2,1.8) -- (0,3.5) -- (3,4.2);
\usetikzlibrary{arrows}
\usetikzlibrary{shapes}
\tikzstyle{every node}=[draw=black,fill=black, ellipse]
\node[minimum width=2.6cm, minimum height = 1cm] (a) at (0,0) {};
\node[minimum width=2cm, minimum height = .8cm] (b) at (4,1.6) {};
\node[minimum width=1.6cm, minimum height = .6cm] (c) at (0,3.2) {};
\node[minimum width=1.2cm, minimum height = .4cm] (d) at (3.2,3.9) {};
\draw[fill = black] (a.75) -- (b.188) -- (b.235) -- (a.06);
\draw[fill = black] (b.170) -- (c.318) -- (c.348) -- (b.140);
\draw[fill = black] (c.5) -- (d.205) -- (d.182) -- (c.28);
\draw[fill = black]  (2.8,5.5) to  [bend right=14] (3.2,4.6) to [bend right =14] (3.6,5.5) -- cycle;
\draw[fill = black] (3.6,5.5) arc (0:180:.4cm);
\draw[fill = white] (3.2,5.5) circle (.22cm);

\end{tikzpicture}



\end{column}
\end{columns}
\end{frame}


\addtobeamertemplate{navigation symbols}{}{%
    \usebeamerfont{footline}%
    \hspace{5em}%
    {\color{black!50}{\small\insertframenumber/\inserttotalframenumber}}%
}
%------------------------------------------------
\section{How Should We Think About Heterogeneous Value Added?}

 % why we care about heterogeneity
\begin{frame}{\textbf{When heterogeneity second-order in fixed effects models...}}


\begin{wideitemize}
    \item Consider the large class of treatment-effect-like estimators based on fixed effects
    \begin{itemize}
        \item Teacher/school value added
        \item Hospital/insurance/provider report cards
        \item Firm and worker productivity
    \end{itemize}
    
    \item But hospitals, firms, and teachers may have comparative advantages
    \begin{itemize}
        \item \textbf{Efficiency:} allocating resources well could make society better off on average
        \item \textbf{Equity:} if social planner also cares about gaps, there may be even more gains
    \end{itemize}
    
    \item There also might be efficiency or distributional implications of various estimation approaches
    
\end{wideitemize}
\end{frame}



\begin{frame}{\textbf{Review: standard VA}}
    \begin{wideitemize}
    \item Assume student test scores, $y_{it}$, are estimable with
    \begin{align*}
        y_{it} &= \gamma_{j(i)} + \beta X{it} +  \varepsilon_{it}
    \end{align*}
    \item $\gamma_{j(i)}$ are fixed-effects estimates for teacher $j$'s VA 
    \begin{itemize}
        \item Typically assume an autocorrelation structure over time
        \item Then regression adjust for shrinkage $\to$ $\tilde{\gamma}_{jt}$
    \end{itemize}
    \item $X_{it}'$ is a vector of controls (conditional independence)
        %lagged score in other subject, lagged average class and school test scores, gender, race, grade, year, English language status, and special education status.
        \begin{itemize}
            \item Sufficiently rich lagged achievement seems to do well \citep{chetty2014measuring1}
        \end{itemize}
    \item $\varepsilon_{it}$ mechanically combines class and student components
    \end{wideitemize}
    
\end{frame}



 % how we think about heterogeneity
\begin{frame}{\textbf{Overview: Assumptions and insights of conceptual framework}}


\begin{wideitemize}
    \item Heterogeneity in value added by {\color<1>{ptb4}\textbf{expected achievement}} is a natural extension
        
    {\tiny \color{gray} \citep{lockwood2009,condie2014teacher,Delgado2020,bates2022teacher}}

    \item<2-> {\color<2>{ptr5}\textbf{We care:}} Two scenarios sufficient to not care about estimating heterogeneity
      \begin{itemize}
        \item If teachers are homogeneous (no surprise)
        \item If classes are homogeneous \textit{and} social preferences utilitarian

        
    \end{itemize}
    \item<3-> {\color<3>{ptr5}\textbf{Class composition:}} Heterogeneity $\implies$ standard VA not composition neutral
      \begin{itemize}
        \item Better-allocated teachers will score higher 
        \item Teachers with certain students may be penalized (akin to model misspecification)
                
    \end{itemize}
    
    \item<4-> {\color<4>{ptr5}\textbf{How to aggregate:}} Heterogeneous effects aggregated with \textbf{explicit} welfare weights
      \begin{itemize}
        \item  Welfare theory allows us to compare reallocations, check regressivity, etc.
        
    \end{itemize} 
\end{wideitemize}


        \hyperlink{ex1}{\beamerbutton{Example One}}
        \hyperlink{ex2}{\beamerbutton{Example Two}}   
\end{frame}



 % how to estimate heterogeneity? here or later
%%%%%%%%%%%%%%%%%%%%%%%%%%%%%%%%%%%%%%%%%%%%%%%%%%%%%%%%

\begin{frame}{\textbf{Measurement if we care about more than the mean}}

\begin{wideitemize}
    \item Two approaches for estimating heterogeneity along the achievement distribution
    
    \item<2-> Option 1: Split students into $k$ bins and estimate $\gamma_{j,k}$ 
    \begin{itemize}
        \item Works great for small number of bins....
        \item But we don't have infinite data, or actually know the \textit{ex ante} expected score
    \end{itemize}

    \item<3-> Option 2: Predict $\hat{y}$ and use kernel estimator:  $y_{i,t}=g_j(\hat{y}_{i,t})$
    \begin{itemize}
        \item Since $\hat{y}_{it} = \hat{\beta} X_{it}$  predicts expected score consistently but with error
        \item Jacknifed following the endogenous stratification literature \citep{abadie2018endogenous}
        \item Much more information, but less precise
    \end{itemize}

\end{wideitemize}

\end{frame}


%------------------------------------------------
\section{Is there Meaningful Heterogeneity in Teacher Effects?}

%------------------------------------------------

%------------------------------------------------
\begin{frame}{\textbf{We use data from teachers in San Diego, California}}


\begin{wideitemize}
    \item Administrative data on universe students in San Diego Unified School District

    \item<2-> Main Sample: 2,165 teachers teaching grades 3-5
        
    \begin{itemize}
        \item Restrict to school years between spring 2003 and 2013 (long term effects)
        \item Require that students have test scores for consecutive years (to estimate VA)
        \item Require that a teacher teaches at least 50 such students (to power heterogeneity)
    \end{itemize}    
        \beamerbutton{Sample Selection Criteria} \beamerbutton{Summary Statistics}
        
        %\hyperlink{data}{\beamerbutton{Data Cleaning}}

    \item<3-> SDUSD has rich data on many variables of interests
        \begin{itemize}
        \item Math and ELA scores, standardized to the mean and variance of California
        \item Long term outcomes: graduation, college enrollment, degree completion
        \item Controls for student characteristics and lagged student, class, and school achievement
    \end{itemize}  
\end{wideitemize}
\end{frame}



%%%%%%%%%%%%%%%%%%%%%%%%%%%%%%%%%%%%%%%%%%%%%%%%%%%%%%%%




\begin{frame}[label=het1,c]{\textbf{Teachers create very heterogeneous value added}}



\begin{center}
\resizebox{.3\textwidth}{!}{$y_{i,j} = \gamma_{j,k(i)} \mathds{1}(i\in j) + \beta_1 X_i + \epsilon_{i,j}$}
\includegraphics[width=.85\textwidth]{slides/slides_pffls/fig1_heterogeneity.pdf}
\end{center}

\hyperlink{example}{\beamerbutton{Hypothesis Tests}}
\end{frame}
%%%%%%%%%%%%%%%%%%%%%%%%%%%%%%%%%%%%%%%%%%%%%%%%%%%
\begin{frame}[c]{\textbf{Long Term Effects}}
\begin{center}

\resizebox{!}{.02\textwidth}{\alt<1>{$y_{i,j} = \tau \hat{\gamma}_{j} + \beta_2 X_i + \nu_{i,j}$}{
$y_{i,j} = \sum_{k_j,k_i} \tau_{k_j,k_i} \hat{\gamma}_{j,k_j}\mathds{1}(k(i) = k_i) + \beta_2 X_i + \nu_{i,j}$}}


\includegraphics<1-2>[width=.85\textwidth]{slides/slides_pffls/fig2a_longterm.pdf}%
\includegraphics<3>[width=.85\textwidth]{slides/slides_pffls/fig2b_longterm.pdf}
\end{center}

\end{frame}



%------------------------------------------------
\section{What are the Implications for Welfare?}


% can i get this to one slide?
%%%%%%%%%%%%%%%%%%%%%%%%%%%%%%%%%%%%%%%%%%%%%%%%%%%%%%%%

\begin{frame}{\textbf{Connecting value added to social welfare}}

\begin{itemize}
    \item We can think of the following social welfare function
    \begin{itemize}
        \item $\omega(x_i,y_i)$ weights: likely based on \textit{ex ante} expected performance
        % Talk about the concavity here?
        \item $v(x_i,y_i)$ value: think achievement, gains, etc.
    \end{itemize}
    \[
    W  = \sum_i \omega(x_i,y_i) v(x_i,y_i) 
    \] 
    
    \item Traditional VAM take the average gains for students among each teacher
    \begin{itemize}
        \item Let $x_i$ be \textit{ex ante} expected performance, estimated as $\hat{y}_i$, then $\tilde{v}(\cdot) = y_i - \hat{y}_i$
    \end{itemize}
    \[
    \hat{W}_{VA}  = \sum_i \frac{(y_i-\hat{y}_i)}{N_{j(i)}} \hspace{3em}
    \]
    
    \item This implies $\omega(x_i,y_i)=\frac{1}{N_j}$, which is (almost) utilitarian

\end{itemize}


\end{frame}

%%%%%%%%%%%%%%%%%%%%%%%%%%%%%%%%%%%%%%%%%%%%%%%%%%%%%%%%

\begin{frame}{\textbf{``Welfare added''}}

\begin{wideitemize}

    \item Let $\hat{v}_j(x_i,y_i)$ estimate a teacher's heterogeneous value added, $\mathbb{E}[v(\cdot)|x_i,j]$
    
    \item Then the estimated ``Welfare Added'' by this teacher is  
    \[
    \hat{W}_j  = \sum_{i\in j} \omega(x_i,y_i) \hat{v}_j(x_i,y_i) 
    \] 
    
    \item Consider the following examples:
    \begin{itemize}
        \item Utilitarian: All students are weighted equally $\omega(x_i,y_i) = 1$
        \item Rawlsian: We care only about students below some cutoff, $c$, $\omega(x_i,y_i) = \mathds{1}(x_i\leq c)$
        \item Beta: Flexible with monotone derivative $ \omega(x_i,y_i)=  \mathcal{B}(x_i;\alpha,\beta)$ 
        
\beamerbutton{Beta Distribution}
    \end{itemize}
    
    %\item But estimating Welfare Added requires an estimate of $\mathbb{E}[v(\cdot)|x_i,j]$

\end{wideitemize}


\end{frame}



\begin{frame}{\textbf{We propose three measures to quantify ``regressivity''?}}

\begin{wideitemize}
\item \textbf{Allocative Regressivity:} Who would benefit from VA-inspired allocations?
\begin{itemize}
    \item Solve optimal allocations as mixed-integer linear programming problem
    \item Compare optimal allocations using traditional and binned VA measures
    \item .... Can we say anything about the actual allocation?
    
\end{itemize}

\item \textbf{Implicit Weights:} How utilitarian are the implicit weights?
\begin{itemize}
    \item Search over welfare functions (beta distribution) to see what yields similar results
    
    \hyperlink{beta}{\beamerbutton{Beta Distribution Examples}}
\end{itemize}

\item \textbf{Teacher Ranks:} Heterogeneity $\implies$ composition affects value added scores
\begin{itemize}
    \item Compare standard value added with composition-invariant utilitarian measures
    \item Simulate wage tax/subsidy under a common performance pay system %\citep{hudson2010effects}
    %teacher advancement program
    \item Compare implications of measurement by race
    
    \hyperlink{tap}{\beamerbutton{Teacher Advancement Program}}
\end{itemize}



\end{wideitemize}

\end{frame}


\begin{frame}[c]{\textbf{Efficiency and equity gains from reallocations}}
\centering
\resizebox{.3\textwidth}{!}{$\max_\mathcal{J} \mathcal{W}(\mathcal{J}) =  \sum_j \sum_{i\in j} \omega_{k(i)} \gamma_{j,{k(i)}}$}
\includegraphics<1>[width=.85\textwidth]{slides/slides_pffls/fig3o_reallocation.pdf}
\includegraphics<2>[width=.85\textwidth]{slides/slides_pffls/fig3a_reallocation.pdf}
\includegraphics<3>[width=.85\textwidth]{slides/slides_pffls/fig3b_reallocation.pdf}
\includegraphics<4>[width=.85\textwidth]{slides/slides_pffls/fig3c_reallocation.pdf}
\includegraphics<5>[width=.85\textwidth]{slides/slides_pffls/fig3d_reallocation.pdf}



\end{frame}


%\begin{frame}{\textbf{Long-Term Gains}}

%\end{frame}




\begin{frame}[c]{\textbf{(Econometric) Implicit weights are not quite utilitarian}}

\centering
\includegraphics<2>[width=.9\textwidth]{slides/slides_pffls/fig6_implicitweights.pdf}

\end{frame}

%similar .55

\begin{frame}[c]{\textbf{Standard measures penalize minority teachers}}

\centering
\includegraphics<2>[width=.9\textwidth]{slides/slides_pffls/fig5_racial.pdf}
% This would raise performance-pay wages by 7% 
\end{frame}


\section{Policy Implications and Conclusion}

\begin{frame}{\textbf{So what did we learn?}}

\begin{wideitemize}
\item Good reason to believe heterogeneity (and estimating heterogeneity) matters
\item<2-> Given heterogeneity, standard VA may not be ``utilitarian'' for students or teachers
\begin{itemize}
    \item Econometrically very close to utilitarian for students
    \item<3-> But penalizes black teachers (because class composition)
\end{itemize}
\item<4-> Teacher effectiveness varies with student achievement---and has long-term implications
\item<5-> Lots of room for gains through reallocations---far beyond standard VA
\end{wideitemize}

\end{frame}

% --------------------------------
\begin{frame}{\textbf{Next steps: Things we could add}}
\begin{wideitemize}

    \item Specification checks vs \citep{chetty2014measuring1,Delgado2020,pope2017multidimensional}
    \item Quasi-experimental validation and placebo tests
    \item What do allocative decisions suggest about social preferences?
    % Hard b/c not on the frontier and we don't know the dimensionality.
    \item Effects beyond of education: teacher wages (unions/not), housing markets?
    \item How do school administrators allocate students to classes?
    %(revealed preference for short gains, long gains, behavior etc?)
    

\end{wideitemize}
   

\end{frame}




\setbeamertemplate{navigation symbols}{} 
\begin{frame}[c,noframenumbering]
\centering
\Huge{\centerline{Thank You!}}

\normalsize  teastmond@ucsd.edu \hspace{1em} njmather@umich.edu \hspace{1em} ricksmi@umich.edu \hspace{1em} jbetts@ucsd.edu

\end{frame}

%------------------------------------------------


\begin{frame}
\frametitle{References}
\tiny
\bibliography{citations}
\end{frame}


%----------------------------------------------------------------------------------------
%----------------------------------------------------------------------------------------

\appendix
\setbeamertemplate{navigation symbols}{}


\section{Bonus Slides}
%----------------------------------------------------------------------------------------
%----------------------------------------------------------------------------------------


\begin{frame}[label=example]{Within Teacher Variation in Value Added}

    \centering
    $H_0: \text{Low Bin VA} = \text{High Bin VA}$
    \resizebox{.8\textwidth}{!}{\renewcommand{\arraystretch}{1.2}

\resizebox{\textwidth}{!}{%
\begin{tabular}{lcc}
    \hline
    \hline
     & $p<0.01$ & $p>0.01$ \\
    \hline
    \hline 
    English Language & 2,015 & 150  \\
    Mathematics & 2,032 & 133 \\
    \hline
    \hline
\end{tabular}}}

\hyperlink{het1}{\beamerbutton{Back}}
\end{frame}



%%%%%%%%%%%%%%%%%%%%%%%%%%%%%%%%%%%%%%%%%%%%%%%%%%%%%%%%
%%%%%%%%%%%%%%%%%%%%%%%%%%%%%%%%%%%%%%%%%%%%%%%%%%%%%%%%




%%%%%%%%%%%%%%%%%%%%%%%%%%%%%%%%%%%%%%%%%%%%%%%%%%%%%%%%




\end{document}



